\documentclass[11pt]{article}
\usepackage{mathblatt}
\def\mitzone{1}
\begin{document}
\blattkopf{Zuordnungen}{Gesamt}
\weit
\verzeichniszeile{\verz{zone}{Kennst du schon (Nr.~1–7)} \verztrenn \verz{e1}{1 Zuordnungen darstellen (Nr.~8–11)} \verztrenn \verz{e2}{2 Proportionale Zuordnungen und Dreisatz (Nr.~12–19)} \verztrenn \verz{e3}{3 Antiproportionale Zuordnungen (Nr.~20–22)} \verztrenn \verz{e4}{4 Zuordnungstypen erkennen und anwenden (Nr.~23–27)} \verztrenn \verz{abhaken}{Das kann ich}}
% Zone „Das kennst du schon“ – aus bank/zuordnungen/zone.jsonl (zusammenbau v0.1)
\einheitenkopf[zone][Kennst du schon]{Das kennst du schon}
\setcounter{aufgabe}{0}

%% TODO Titel: Ich-kann-Satz fehlt in der Bank (Platzhalter: Kettenname) – Nr. 1
%% TODO Anweisung: ein Satz über der ersten Teilaufgabe fehlt in der Bank – Nr. 1
\begin{aufgabe}{Punkte im Koordinatensystem eintragen und ablesen (erster Quadrant)}
\begin{teile}
\teil Trage den Punkt A(4|3) ins Koordinatensystem ein.

\begin{ksys}[xmin=0,xmax=6,ymin=0,ymax=6]\end{ksys}
\teil Lies ab: Welche Koordinaten hat der Punkt C?

\begin{ksys}[xmin=0,xmax=6,ymin=0,ymax=6,ablesen]\punkt{2.5}{4}{C}\end{ksys}
C( \leerfeld | \leerfeld )
\end{teile}
\end{aufgabe}

%% TODO Titel: Ich-kann-Satz fehlt in der Bank (Platzhalter: Kettenname) – Nr. 2
%% TODO Anweisung: ein Satz über der ersten Teilaufgabe fehlt in der Bank – Nr. 2
\begin{aufgabe}{Multiplizieren und Dividieren mit Dezimalzahlen (Preis geteilt durch Stückzahl, Preis mal Stückzahl)}
\begin{teile}
\teil 3 Hefte kosten je 1,50 €. Was kosten sie zusammen? \leerfeld[€]
\teil 7 Kugeln Eis kosten je 1,20 €. Was kosten alle zusammen? \leerfeld[€]
\end{teile}
\end{aufgabe}

%% TODO Titel: Ich-kann-Satz fehlt in der Bank (Platzhalter: Kettenname) – Nr. 3
%% TODO Anweisung: ein Satz über der ersten Teilaufgabe fehlt in der Bank – Nr. 3
\begin{aufgabe}{Vielfache und Teiler erkennen (zwölf ist das Dreifache von vier)}
\begin{teile}
\teil 15 ist das Wievielfache von 5? das \leerfeld -Fache
\teil 3 kg sind das Wievielfache von 0,5 kg? das \leerfeld -Fache
\end{teile}
\end{aufgabe}

%% TODO Titel: Ich-kann-Satz fehlt in der Bank (Platzhalter: Kettenname) – Nr. 4
%% TODO Anweisung: ein Satz über der ersten Teilaufgabe fehlt in der Bank – Nr. 4
\begin{aufgabe}{Bruchteil einer Größe (die Hälfte, ein Viertel von 12 €)}
\begin{teile}
\teil Die Hälfte von 14 € – wie viel? \leerfeld[€]
\teil Ein Viertel von 3 € – wie viel? \leerfeld[€]
\end{teile}
\end{aufgabe}

%% TODO Titel: Ich-kann-Satz fehlt in der Bank (Platzhalter: Kettenname) – Nr. 5
%% TODO Anweisung: ein Satz über der ersten Teilaufgabe fehlt in der Bank – Nr. 5
\begin{aufgabe}{Einheiten umrechnen (Cent ↔ Euro, Minuten ↔ Stunden, m ↔ km)}
\begin{teile}
\teil Wie viele Minuten sind 3 Stunden? \leerfeld[min]
\teil Wie viele Meter sind 4,2 km? \leerfeld[m]
\end{teile}
\end{aufgabe}

%% TODO Titel: Ich-kann-Satz fehlt in der Bank (Platzhalter: Kettenname) – Nr. 6
%% TODO Anweisung: ein Satz über der ersten Teilaufgabe fehlt in der Bank – Nr. 6
\begin{aufgabe}{Vielfache und Teiler erkennen (zwölf ist das Dreifache von vier) – Fehler finden}
\begin{teile}
\teil Finn schreibt: „Von 7 € auf 28 € sind es 21 € mehr, also sind 28 € das 21-Fache von 7 €.“ Finde den Fehler und rechne richtig.
\end{teile}
\end{aufgabe}

%% TODO Titel: Ich-kann-Satz fehlt in der Bank (Platzhalter: Kettenname) – Nr. 7
%% TODO Anweisung: ein Satz über der ersten Teilaufgabe fehlt in der Bank – Nr. 7
\begin{aufgabe}{Vielfache und Teiler erkennen (zwölf ist das Dreifache von vier) – selbst rechnen}
\begin{teile}
\teil Ein Beet war 9 m² groß, jetzt ist es 45 m² groß. Das Wievielfache ist das? das \leerfeld -Fache
\end{teile}
\end{aufgabe}

\clearpage
% Einheit 1 von 4 – Zuordnungen darstellen (bank/zuordnungen/e1.jsonl, zusammenbau v0.1)
%% TODO Zweigzeile Teil 1: was hier gelernt wird (Satz in Schülersprache aus dem Einheitstitel) – Einheit 1
\einheitenkopf[e1]{Einheit 1 von 4 · Zuordnungen darstellen}
\zweigzeile{ab Kl. 6 · P10}
\setcounter{aufgabe}{7}

%% TODO \verfahren{…}: Verfahrensüberschrift in Sachsprache fehlt in der Bank (2.3 b)
%% TODO Titel: Ich-kann-Satz fehlt in der Bank (Platzhalter: Kettenname) – Nr. 8
%% TODO Anweisung: ein Satz über der ersten Teilaufgabe fehlt in der Bank – Nr. 8
\begin{aufgabe}{Darstellen}
%% TODO \rechenplatz{n}: Schrittzahl des Grundfalls fehlt in der Bank (nur bei fester Schreibform, 2.3 b)
\begin{teile}
\teil Ein Radfahrer fährt gleichmäßig: nach 1 Stunde 15 km, nach 2 Stunden 30 km, nach 3 Stunden 45 km. Trage die Wertepaare in die Tabelle ein.

\wertetabelleleer{Zeit in h}{Weg in km}{3}
\teil Trage die Punkte aus der Tabelle ins Koordinatensystem ein.

\wertetabelle[2,4,6,8]{Zeit in h}{Weg in km}{1,2,3,4} \begin{ksys}[xmin=0,xmax=5,ymin=0,ymax=9,xlabel=Zeit in h,ylabel=Weg in km]\end{ksys}
\teil Lies ab: Wie weit ist man nach 3 Stunden? Nach wie vielen Stunden sind es 8 km?

\begin{ksys}[xmin=0,xmax=5,ymin=0,ymax=10,ablesen,xlabel=Zeit in h,ylabel=Weg in km]\gerade{2}{0}{}\end{ksys}
\leerfeld[km] ; \leerfeld[h]
\teil Trage die Punkte aus der Tabelle ins Koordinatensystem ein. Senkrecht steht ein Kästchen für 10 l.

\wertetabelle[20,40,60,80]{Zeit in min}{Wasser in l}{1,2,3,4} \begin{ksys}[xmin=0,xmax=5,ymin=0,ymax=100,ystep=10,xlabel=Zeit in min,ylabel=Wasser in l]\end{ksys}
\teil Die Punkte gehören zur Zuordnung Anzahl Kinokarten → Preis. Dürfen sie zu einer Linie verbunden werden? Begründe.
\teil Ein Eimer wird mit einem gleichmäßigen Wasserstrahl gefüllt. Welcher Graph passt zu Zeit → Füllhöhe? \\ \kreuz{steigende Gerade durch den Ursprung} \\ \kreuz{fallende Gerade} \\ \kreuz{waagerechte Linie}
\teil Mia hat die Tabelle ins Koordinatensystem übertragen (Bild). Finde den Fehler.

\wertetabelle[3,6]{Zeit in h}{Weg in km}{1,2} \begin{ksys}[xmin=0,xmax=7,ymin=0,ymax=7,xlabel=Zeit in h,ylabel=Weg in km]\punkt{3}{1}{A}\punkt{6}{2}{B}\end{ksys}
\teil Ein Sportverein zählt die Abonnenten seines Kanals. Trage die Werte aus der Tabelle ins Koordinatensystem ein. Die senkrechte Achse ist in Schritten von 50 beschriftet, ein Kästchen steht für 10. \hfill (P10 2025 OS)

\wertetabelle[70,95,135,185,260]{Zeit in Monaten}{Abonnenten}{0,1,2,3,4} \begin{ksys}[xmin=0,xmax=5,ymin=0,ymax=300,ystep=50,xlabel=Zeit in Monaten,ylabel=Abonnenten]\end{ksys}
\end{teile}
\end{aufgabe}

%% TODO \verfahren{…}: Verfahrensüberschrift in Sachsprache fehlt in der Bank (2.3 b)
%% TODO Titel: Ich-kann-Satz fehlt in der Bank (Platzhalter: Kettenname) – Nr. 9
%% TODO Anweisung: ein Satz über der ersten Teilaufgabe fehlt in der Bank – Nr. 9
\begin{aufgabe}{Achseneinteilung}
%% TODO \rechenplatz{n}: Schrittzahl des Grundfalls fehlt in der Bank (nur bei fester Schreibform, 2.3 b)
\begin{teile}
\teil Auf einer Achse stehen die Zahlen 0, 20 und 40, je 4 Kästchen auseinander. Wie viel ist ein Kästchen wert? 1 Kästchen = \leerfeld
\teil Die Werte gehen bis 60 km, die Achse hat 15 Kästchen. Welche Einteilung passt? \\ \kreuz{1 Kästchen = 1 km} \\ \kreuz{1 Kästchen = 5 km}
\teil Wähle für die waagerechte Achse eine Einteilung und zeichne die beschriftete Achse.

\wertetabelle[2,4,6,8,10,12]{Anzahl Hefte}{Preis in €}{1,2,3,4,5,6}
\teil Der größte Wert ist 90 km, ein Kästchen soll 5 km sein. Wie viele Kästchen braucht man? Passt es auf eine Achse mit 16 Kästchen? \leerfeld[Kästchen]
\teil Die Werte gehen bis 800 m, die Achse hat 20 Kästchen. Wie viele Meter muss ein Kästchen mindestens tragen? 1 Kästchen = \leerfeld[m]
\teil Zeichne die beiden Achsen für diese Tabelle und beschrifte sie mit Größe und Einheit.

\wertetabelle[400,800,1\,200]{Zeit in min}{Strecke in m}{5,10,15}
\teil Ein Regenfass enthält anfangs 90 l Wasser und läuft gleichmäßig aus. Ergänze die Tabelle. Zeichne dann auf Karopapier ein Koordinatensystem mit 20 Kästchen nach rechts und 20 Kästchen nach oben: Wähle eine Einteilung, beschrifte beide Achsen und trage die Punkte ein. \hfill (P10 2021 OS)

\wertetabelle[,80,70,60,50,40,]{Zeit in min}{Inhalt in l}{0,5,10,15,20,25,30}
\end{teile}
\end{aufgabe}

%% TODO Titel: Ich-kann-Satz fehlt in der Bank (Platzhalter: Kettenname) – Nr. 10
%% TODO Anweisung: ein Satz über der ersten Teilaufgabe fehlt in der Bank – Nr. 10
\begin{aufgabe}{Formel aus Tabelle (y = 4 · x)}
\begin{teile}
\teil Die Tabelle gehört zu einer proportionalen Zuordnung (x: Anzahl Hefte, y: Preis in €). Wie lautet die Formel?

\wertetabelle[3,6,9]{x}{y}{1,2,3}
y = \leerfeld · x
\end{teile}
\end{aufgabe}

%% TODO Titel: Ich-kann-Satz fehlt in der Bank (Platzhalter: Kettenname) – Nr. 11
%% TODO Anweisung: ein Satz über der ersten Teilaufgabe fehlt in der Bank – Nr. 11
\begin{aufgabe}{Darstellen – Fehler finden, Begründen, Darstellungswechsel, Anwendung}
\begin{teile}
\teil Lea liest am Graphen ab, wie weit sie nach 2 Stunden ist. Der Punkt liegt 5 Kästchen hoch, senkrecht steht ein Kästchen für 2 km. Lea schreibt: 5 km. Finde den Fehler.
\teil Kann in einer Tabelle Uhrzeit → Temperatur dieselbe Uhrzeit zweimal mit verschiedenen Temperaturen stehen? Begründe.
\teil Lies am Graphen die Werte für 1, 2 und 3 Stunden ab und schreibe sie als Tabelle.

\begin{ksys}[xmin=0,xmax=4,ymin=0,ymax=6,ablesen,xlabel=Zeit in h,ylabel=Regen in mm]\gerade{1.5}{0}{}\end{ksys}
\teil Das Diagramm zeigt die Temperatur an einem Sommertag. Lies ab: Um wie viel Uhr war es am wärmsten? Um wie viel Grad stieg die Temperatur von 8 Uhr bis dahin?

\liniendia[ymax=25,ystep=5,yfein=1,ylabel=Temperatur in $^\circ$C]{8 Uhr/12,10 Uhr/16,12 Uhr/20,14 Uhr/23,16 Uhr/21,18 Uhr/17}
\leerfeld Uhr; \leerfeld °C
\end{teile}
\end{aufgabe}

\clearpage
% Einheit 2 von 4 – Proportionale Zuordnungen und Dreisatz (bank/zuordnungen/e2.jsonl, zusammenbau v0.1)
%% TODO Zweigzeile Teil 1: was hier gelernt wird (Satz in Schülersprache aus dem Einheitstitel) – Einheit 2
\einheitenkopf[e2]{Einheit 2 von 4 · Proportionale Zuordnungen und Dreisatz}
\zweigzeile{ab Kl. 6 · P10}
\setcounter{aufgabe}{11}

%% TODO Titel: Ich-kann-Satz fehlt in der Bank (Platzhalter: Kettenname) – Nr. 12
%% TODO Anweisung: ein Satz über der ersten Teilaufgabe fehlt in der Bank – Nr. 12
\begin{aufgabe}{Was ist hier eine Portion?}
\begin{teile}
\teil Äpfel: 3 kg kosten 4,50 €. Was ist hier eine Portion?
\end{teile}
\end{aufgabe}

%% TODO Titel: Ich-kann-Satz fehlt in der Bank (Platzhalter: Kettenname) – Nr. 13
%% TODO Anweisung: ein Satz über der ersten Teilaufgabe fehlt in der Bank – Nr. 13
\begin{aufgabe}{Minitabelle anlegen}
\begin{teile}
\teil 5 kg Kartoffeln kosten 6 €. Was kosten 8 kg? Rechne noch nicht: Welche Mengen stehen in der Minitabelle links untereinander?
\end{teile}
\end{aufgabe}

%% TODO Titel: Ich-kann-Satz fehlt in der Bank (Platzhalter: Kettenname) – Nr. 14
%% TODO Anweisung: ein Satz über der ersten Teilaufgabe fehlt in der Bank – Nr. 14
\begin{aufgabe}{Pfeile beidseitig}
\begin{teile}
\teil 8 Hefte kosten 10 €. Im Dreisatz steht links am Pfeil „: 8“. Was steht rechts am Pfeil?
\end{teile}
\end{aufgabe}

%% TODO Titel: Ich-kann-Satz fehlt in der Bank (Platzhalter: Kettenname) – Nr. 15
%% TODO Anweisung: ein Satz über der ersten Teilaufgabe fehlt in der Bank – Nr. 15
\begin{aufgabe}{Dreisatz proportional}
%% TODO \rechenplatz{n}: Schrittzahl des Grundfalls fehlt in der Bank (nur bei fester Schreibform, 2.3 b)
\begin{teile}
\teil Führe die Tabelle weiter.

\wertetabelle[2,4,,]{Anzahl Stück}{Preis in €}{1,2,3,4}
\teil 2 kg Birnen kosten 3 €. Was kosten 4 kg? \leerfeld[€]
\teil 5 Joghurts kosten 3 €. Was kosten 15 Joghurts? \leerfeld[€]
\teil 30 Hefte kosten 45 €. Was kosten 10 Hefte? \leerfeld[€]
\teil 3 Kugeln Eis kosten 6 €. Was kosten 7 Kugeln? \leerfeld[€]
\teil 5 Brezeln kosten 3,50 €. Was kosten 9 Brezeln? \leerfeld[€]
\teil Packung A: 500 g Nudeln für 2,00 €. Packung B: 750 g für 2,70 €. Welche ist günstiger? \\ \kreuz{Packung A} \\ \kreuz{Packung B}
\teil Ein Läufer braucht für 2 km 13 Minuten. Wie lange braucht er bei gleichem Tempo für 5 km? Rechne ohne Taschenrechner. \leerfeld[min]
\end{teile}
\end{aufgabe}

%% TODO Titel: Ich-kann-Satz fehlt in der Bank (Platzhalter: Kettenname) – Nr. 16
%% TODO Anweisung: ein Satz über der ersten Teilaufgabe fehlt in der Bank – Nr. 16
\begin{aufgabe}{fester Faktor angeben (Preis je Einheit), Gleichung y = k · x}
\begin{teile}
\teil 7 kg Äpfel kosten 17,50 €. Gib den Preis je kg an und schreibe die Gleichung (x: Menge in kg, y: Preis in €). \leerfeld[€] je kg; y = \leerfeld · x
\end{teile}
\end{aufgabe}

%% TODO Titel: Ich-kann-Satz fehlt in der Bank (Platzhalter: Kettenname) – Nr. 17
%% TODO Anweisung: ein Satz über der ersten Teilaufgabe fehlt in der Bank – Nr. 17
\begin{aufgabe}{Graph zeichnen (Ursprungsgerade) und Werte ablesen}
\begin{teile}
\teil 1 kg Kirschen kostet 3 €. Zeichne den Graphen Menge in kg → Preis in € und lies ab, was 2,5 kg kosten.

\begin{ksys}[xmin=0,xmax=5,ymin=0,ymax=15,xlabel=Menge in kg,ylabel=Preis in €]\end{ksys}
\end{teile}
\end{aufgabe}

%% TODO Titel: Ich-kann-Satz fehlt in der Bank (Platzhalter: Kettenname) – Nr. 18
%% TODO Anweisung: ein Satz über der ersten Teilaufgabe fehlt in der Bank – Nr. 18
\begin{aufgabe}{Verhältnisgleichung aufstellen (Vorrat)}
\begin{gleichungsraster}[2]
\gl{\text{Für 4 Portionen Pfannkuchenteig braucht man 300 g Mehl. Wie viel Mehl braucht man für 10 Portionen? Stelle eine Verhältnisgleichung auf und löse sie (x: Mehl in g).}} & \\
\end{gleichungsraster}
\end{aufgabe}

%% TODO Titel: Ich-kann-Satz fehlt in der Bank (Platzhalter: Kettenname) – Nr. 19
%% TODO Anweisung: ein Satz über der ersten Teilaufgabe fehlt in der Bank – Nr. 19
\begin{aufgabe}{Dreisatz proportional – Fehler finden, Begründen, Darstellungswechsel, Anwendung}
\begin{teile}
\teil 3 kg Kartoffeln kosten 7,20 €. Tim rechnet für 1 kg: 7,20 € − 3 = 4,20 €. Finde den Fehler und rechne richtig.
\teil Warum rechnet man im Dreisatz rechts auch „: 4“, wenn man links „: 4“ rechnet? Begründe.
\teil Der Graph zeigt Menge in kg → Preis in €. Lies den Preis für 1 kg ab und schreibe die Gleichung (x: Menge in kg, y: Preis in €).

\begin{ksys}[xmin=0,xmax=4,ymin=0,ymax=6,ablesen,xlabel=Menge in kg,ylabel=Preis in €]\gerade{1.5}{0}{}\end{ksys}
\leerfeld[€] je kg; y = \leerfeld · x
\teil Ein Rezept für 4 Personen braucht 250 g Nudeln. Zum Geburtstag kommen 10 Personen. Reicht eine Packung mit 500 g?
\end{teile}
\end{aufgabe}

\clearpage
% Einheit 3 von 4 – Antiproportionale Zuordnungen (bank/zuordnungen/e3.jsonl, zusammenbau v0.1)
%% TODO Zweigzeile Teil 1: was hier gelernt wird (Satz in Schülersprache aus dem Einheitstitel) – Einheit 3
\einheitenkopf[e3]{Einheit 3 von 4 · Antiproportionale Zuordnungen}
\zweigzeile{ab Kl. 7 · P10}
\setcounter{aufgabe}{19}

%% TODO Titel: Ich-kann-Satz fehlt in der Bank (Platzhalter: Kettenname) – Nr. 20
%% TODO Anweisung: ein Satz über der ersten Teilaufgabe fehlt in der Bank – Nr. 20
\begin{aufgabe}{Antiproportional}
%% TODO \rechenplatz{n}: Schrittzahl des Grundfalls fehlt in der Bank (nur bei fester Schreibform, 2.3 b)
\begin{teile}
\teil Je mehr Maler an einer Wand arbeiten, desto … Zeit brauchen sie. \\ \kreuz{mehr} \\ \kreuz{weniger}
\teil 2 Bagger brauchen für eine Baugrube 10 Tage. Wie lange brauchen 4 Bagger? \leerfeld[Tage]
\teil 2 Traktoren pflügen ein Feld in 9 Stunden. Wie lange brauchen 6 Traktoren? \leerfeld[h]
\teil 3 Maschinen brauchen für einen Auftrag 8 Stunden. Wie lange brauchen 4 Maschinen? \leerfeld[h]
\teil Jemand behauptet: 6 Arbeiter brauchen 10 Tage, 4 Arbeiter brauchen 15 Tage. Prüfe mit dem Produkt, ob das zusammenpasst.
\teil Für einen Umzug haben 3 Freunde 4 Stunden eingeplant. Kurzfristig helfen 3 weitere Freunde mit. Alle arbeiten gleich schnell. Wie lange dauert der Umzug jetzt?
\teil Ein Futtervorrat reicht für 6 Ziegen 10 Tage. Wie viele Tage reicht er für 5 Ziegen? \hfill (P10 2014 GYM) \\ \leerfeld[Tage]
\end{teile}
\end{aufgabe}

%% TODO Titel: Ich-kann-Satz fehlt in der Bank (Platzhalter: Kettenname) – Nr. 21
%% TODO Anweisung: ein Satz über der ersten Teilaufgabe fehlt in der Bank – Nr. 21
\begin{aufgabe}{Graph als fallende Kurve, Werte ablesen}
\begin{teile}
\teil Lies ab: Wie lange brauchen 3 Helfer? Wie viele Helfer braucht man, damit es 2 Stunden dauert?

\begin{ksys}[xmin=0,xmax=7,ymin=0,ymax=13,ablesen,xlabel=Anzahl Helfer,ylabel=Zeit in h]\funktionab{12/\x}{}{1}{6}\end{ksys}
\leerfeld[h] ; \leerfeld Helfer
\end{teile}
\end{aufgabe}

%% TODO Titel: Ich-kann-Satz fehlt in der Bank (Platzhalter: Kettenname) – Nr. 22
%% TODO Anweisung: ein Satz über der ersten Teilaufgabe fehlt in der Bank – Nr. 22
\begin{aufgabe}{Antiproportional – Fehler finden, Begründen, Darstellungswechsel, Anwendung}
\begin{teile}
\teil 4 Maler brauchen 5 Tage. Kim rechnet: „8 Maler brauchen 10 Tage, denn es sind doppelt so viele.“ Finde den Fehler und rechne richtig.
\teil Warum bleibt bei der Zuordnung Anzahl Arbeiter → Anzahl Tage das Produkt gleich? Begründe.
\teil Trage die Punkte der Tabelle ein und verbinde sie zu einer Kurve.

\wertetabelle[6,3,2,1]{Anzahl Kinder}{Stück je Kind}{1,2,3,6} \begin{ksys}[xmin=0,xmax=7,ymin=0,ymax=7,xlabel=Anzahl Kinder,ylabel=Stück je Kind]\end{ksys}
\teil Für das Schmücken der Turnhalle plant die Klasse mit 5 Helfern 3 Stunden. Die Halle ist aber nur 2 Stunden frei. Wie viele Helfer braucht man mindestens?
\end{teile}
\end{aufgabe}

\clearpage
% Einheit 4 von 4 – Zuordnungstypen erkennen und anwenden (bank/zuordnungen/e4.jsonl, zusammenbau v0.1)
%% TODO Zweigzeile Teil 1: was hier gelernt wird (Satz in Schülersprache aus dem Einheitstitel) – Einheit 4
\einheitenkopf[e4]{Einheit 4 von 4 · Zuordnungstypen erkennen und anwenden}
\zweigzeile{ab Kl. 7 · P10 oft}
\setcounter{aufgabe}{22}

%% TODO Titel: Ich-kann-Satz fehlt in der Bank (Platzhalter: Kettenname) – Nr. 23
%% TODO Anweisung: ein Satz über der ersten Teilaufgabe fehlt in der Bank – Nr. 23
\begin{aufgabe}{Nullwert prüfen}
\begin{teile}
\teil Taxi: 4 € Grundgebühr und 2 € je km. Kostet eine Fahrt von 0 km auch 0 €? \\ \kreuz{ja} \\ \kreuz{nein}
\end{teile}
\end{aufgabe}

%% TODO \verfahren{…}: Verfahrensüberschrift in Sachsprache fehlt in der Bank (2.3 b)
%% TODO Titel: Ich-kann-Satz fehlt in der Bank (Platzhalter: Kettenname) – Nr. 24
%% TODO Anweisung: ein Satz über der ersten Teilaufgabe fehlt in der Bank – Nr. 24
\begin{aufgabe}{Erkennen}
%% TODO \rechenplatz{n}: Schrittzahl des Grundfalls fehlt in der Bank (nur bei fester Schreibform, 2.3 b)
\begin{teile}
\teil Je mehr Hefte du kaufst, desto … bezahlst du. \\ \kreuz{mehr} \\ \kreuz{weniger} \\ \kreuz{nicht vorhersagbar}
\teil Ist die Zuordnung in der Tabelle proportional? Prüfe die Quotienten y : x. \\ \kreuz{proportional} \\ \kreuz{nicht proportional}

\wertetabelle[6,12,15]{x}{y}{2,4,5}
\teil Ist die Zuordnung in der Tabelle antiproportional? Prüfe die Produkte x · y. \\ \kreuz{antiproportional} \\ \kreuz{nicht antiproportional}

\wertetabelle[9,6,3]{x}{y}{2,3,6}
\teil Fahrradverleih: 3 € Grundgebühr und 2 € je Stunde. Ist Mietdauer → Preis proportional? \\ \kreuz{proportional} \\ \kreuz{nicht proportional}
\teil Welcher Zuordnungstyp gehört zum Graphen? \\ \kreuz{proportional} \\ \kreuz{antiproportional} \\ \kreuz{keins von beiden}

\begin{ksys}[xmin=0,xmax=5,ymin=0,ymax=8]\gerade{1.5}{0}{}\end{ksys}
\teil Welcher Zuordnungstyp liegt in der Tabelle vor? Prüfe erst die Quotienten, dann die Produkte. \\ \kreuz{proportional} \\ \kreuz{antiproportional} \\ \kreuz{keins von beiden}

\wertetabelle[20,10,5]{x}{y}{2,4,8}
\teil Bei einem Paketdienst kosten 2 kg 8 € und 4 kg 14 €. Ist Gewicht → Preis proportional? Entscheide und begründe.
\end{teile}
\end{aufgabe}

%% TODO \verfahren{…}: Verfahrensüberschrift in Sachsprache fehlt in der Bank (2.3 b)
%% TODO Titel: Ich-kann-Satz fehlt in der Bank (Platzhalter: Kettenname) – Nr. 25
%% TODO Anweisung: ein Satz über der ersten Teilaufgabe fehlt in der Bank – Nr. 25
\begin{aufgabe}{Rate}
%% TODO \rechenplatz{n}: Schrittzahl des Grundfalls fehlt in der Bank (nur bei fester Schreibform, 2.3 b)
\begin{teile}
\teil Auf dem Tankbeleg steht: 40 l, 1,79 € je Liter, gesamt 71,60 €. Welche Angabe ist die Rate?
\teil Äpfel kosten 2,50 € je kg. Was kosten 4 kg? \leerfeld[€]
\teil Kartoffeln kosten 1,50 € je kg. Wie viele Kilogramm bekommt man für 6 €? \leerfeld[kg]
\teil Ein Radfahrer fährt 45 km in 3 Stunden. Wie schnell fährt er im Schnitt? \leerfeld km/h
\teil Ein Bus fährt 150 km mit durchschnittlich 50 km/h. Wie lange dauert die Fahrt? \leerfeld[h]
\teil Ein Lieferwagen verbraucht 8 l Diesel auf 100 km und fährt im Jahr 30\,000 km. Ein Liter Diesel kostet 162,5 ct. Berechne die Dieselkosten für ein Jahr in Euro. \hfill (P10 2014 OS) \\ \leerfeld[€]
\teil Eine Wandergruppe geht mit 4 km/h. Sie ist um 9:50 Uhr an einer Hütte, macht dort 30 Minuten Pause und geht dann 3\,000 m bis zum See. Wann kommt sie am See an? \hfill (P10 2014 OS) \\ um \leerfeld Uhr
\end{teile}
\end{aufgabe}

%% TODO Titel: Ich-kann-Satz fehlt in der Bank (Platzhalter: Kettenname) – Nr. 26
%% TODO Anweisung: ein Satz über der ersten Teilaufgabe fehlt in der Bank – Nr. 26
\begin{aufgabe}{Gegenbeispiele (Alter und Größe, Tarif mit Grundgebühr)}
\begin{teile}
\teil Nenne ein Beispiel für eine Zuordnung „je mehr …, desto mehr …“, die nicht proportional ist, und begründe kurz.
\end{teile}
\end{aufgabe}

%% TODO Titel: Ich-kann-Satz fehlt in der Bank (Platzhalter: Kettenname) – Nr. 27
%% TODO Anweisung: ein Satz über der ersten Teilaufgabe fehlt in der Bank – Nr. 27
\begin{aufgabe}{Erkennen – Fehler finden, Begründen, Darstellungswechsel, Anwendung}
\begin{teile}
\teil Ein Auto fährt 90 km in 1 Stunde. Lukas rechnet: „Für 180 km braucht es eine halbe Stunde, denn die Strecke ist doppelt so lang.“ Finde den Fehler und rechne richtig.
\teil Warum ist die Zuordnung Anzahl Personen → Preis pro Person für eine Pizza nicht proportional? Begründe.
\teil Ein Brunnen liefert 5 l Wasser je Minute. Zeichne den Graphen Zeit in min → Wasser in l für 0 bis 4 Minuten und nenne den Zuordnungstyp.

\begin{ksys}[xmin=0,xmax=5,ymin=0,ymax=20,xlabel=Zeit in min,ylabel=Wasser in l]\end{ksys}
\teil Familie Yilmaz fährt 280 km in den Urlaub. Sie fährt im Schnitt 80 km/h und will um 14:00 Uhr ankommen. Wann muss sie spätestens losfahren?
\end{teile}
\end{aufgabe}

% Abhakseite „Das kann ich“ – Zone nur im Gesamt (\mitzone)
%% TODO Ich-kann-Sätze: Platzhalter wie in den Aufgabendateien
\begin{abhakseite}
\ifdefined\mitzone
\abhakgruppe{Kennst du schon}
\abhak{1}{Punkte im Koordinatensystem eintragen und ablesen (erster Quadrant)}
\abhak{2}{Multiplizieren und Dividieren mit Dezimalzahlen (Preis geteilt durch Stückzahl, Preis mal Stückzahl)}
\abhak{3}{Vielfache und Teiler erkennen (zwölf ist das Dreifache von vier)}
\abhak{4}{Bruchteil einer Größe (die Hälfte, ein Viertel von 12 €)}
\abhak{5}{Einheiten umrechnen (Cent ↔ Euro, Minuten ↔ Stunden, m ↔ km)}
\abhak{6}{Vielfache und Teiler erkennen (zwölf ist das Dreifache von vier) – Fehler finden}
\abhak{7}{Vielfache und Teiler erkennen (zwölf ist das Dreifache von vier) – selbst rechnen}
\fi
\abhakgruppe{Einheit 1 · Zuordnungen darstellen}
\abhak{8}{Darstellen}
\abhak{9}{Achseneinteilung}
\abhak{10}{Formel aus Tabelle (y = 4 · x)}
\abhak{11}{Darstellen – Fehler finden, Begründen, Darstellungswechsel, Anwendung}
\abhakgruppe{Einheit 2 · Proportionale Zuordnungen und Dreisatz}
\abhak{12}{Was ist hier eine Portion?}
\abhak{13}{Minitabelle anlegen}
\abhak{14}{Pfeile beidseitig}
\abhak{15}{Dreisatz proportional}
\abhak{16}{fester Faktor angeben (Preis je Einheit), Gleichung y = k · x}
\abhak{17}{Graph zeichnen (Ursprungsgerade) und Werte ablesen}
\abhak{18}{Verhältnisgleichung aufstellen (Vorrat)}
\abhak{19}{Dreisatz proportional – Fehler finden, Begründen, Darstellungswechsel, Anwendung}
\abhakgruppe{Einheit 3 · Antiproportionale Zuordnungen}
\abhak{20}{Antiproportional}
\abhak{21}{Graph als fallende Kurve, Werte ablesen}
\abhak{22}{Antiproportional – Fehler finden, Begründen, Darstellungswechsel, Anwendung}
\abhakgruppe{Einheit 4 · Zuordnungstypen erkennen und anwenden}
\abhak{23}{Nullwert prüfen}
\abhak{24}{Erkennen}
\abhak{25}{Rate}
\abhak{26}{Gegenbeispiele (Alter und Größe, Tarif mit Grundgebühr)}
\abhak{27}{Erkennen – Fehler finden, Begründen, Darstellungswechsel, Anwendung}
\end{abhakseite}

\end{document}
